\section{Systemarkitektur}

Dette kapitel beskriver Træningsmesters systemarkitektur som teknisk svar på kravene fra det foregående kapitel. Arkitekturen skal gøre det muligt at forbinde programstruktur, faktisk træning, completion, progression og feedback uden at gøre brugerfladen til den endelige kilde til sandhed. Den skal samtidig understøtte lav friktion i træningssituationen og beskytte brugerbundne træningsdata gennem en tydelig backendgrænse.

Arkitekturen behandles derfor som en designmæssig og teknisk struktur, ikke som dokumentation for sundhedseffekt. Dens kvalitet vurderes ud fra, om den skaber sammenhæng mellem krav, dataflow, adgangskontrol og brugerens centrale træningshandlinger.

\subsection{Systemkontekst}

Træningsmester består af en native iOS-app, en watchOS companion, Live Activity- og widgetspor, Supabase som backend samt udvalgte Apple-services. Figur~\ref{fig:tm-system-context} viser denne systemkontekst, mens Tabel~\ref{tab:tm-container-architecture} samler hovedkomponenterne i et tekstligt overblik. Pointen er, at appen ikke kun er en enkelt skærmflade, men et økosystem af klienter, platformservices og backendkomponenter.

iOS-appen er den primære brugerflade for programopbygning, dagens træning, tracker, historik og indstillinger. watchOS-sporet fungerer som en træningsnær companion, hvor brugeren kan få adgang til relevante workoutdata uden nødvendigvis at bruge telefonen aktivt. Live Activity-sporet giver en statusflade under aktiv træning. Supabase håndterer autentifikation, database, Row Level Security, RPC og Storage, mens Edge Functions bruges til privilegerede eller integrationsnære handlinger.

Denne systemkontekst følger kravkapitlets skelnen mellem plan, faktisk udførelse og næste handling. Arkitekturen skal kunne håndtere, at brugeren følger planen, afviger fra den, afslutter en træning uden fuld trackerlog eller bruger flere flader under samme træningsforløb.

\subsection{Lagdeling mellem brugerflade, state og dataadgang}

Arkitekturen er lagdelt for at reducere kobling mellem brugerflade og backend. SwiftUI understøtter deklarative brugerflader og state-drevet rendering \cite{apple_swiftui_docs}, men i Træningsmester er brugerfladen ikke designet som sikkerhedsgrænse eller datamæssig autoritet. Den viser og indsamler handlinger, mens state, domænemodeller, repositories og backendregler afgør, hvad handlingerne betyder.

På konceptuelt niveau kan arkitekturen læses som seks lag: presentation, AppState, domænemodeller, repository-lag, platformintegrationer og backend. Presentation-laget rummer skærme og flows. AppState koordinerer session, brugerprofil, aktiv plan, aktiv træning og centrale runtime-tilstande. Domænemodellerne beskriver de begreber, appen arbejder med, eksempelvis program, workout, øvelse, sæt, completion og træningscyklus. Repository-laget samler dataadgang og oversætter mellem appens domæne og backendens kontrakter.

Denne opdeling er vigtig, fordi træningsflowet ikke bør afhænge af, at en enkelt skærm husker al kontekst. Når workoutstatus, completion og progression er placeret i delte state- og datalag, kan samme træningsforløb understøttes af iOS, watchOS og Live Activity uden at hver flade skal genopfinde logikken.

\subsection{AppState og repository-lag}

AppState fungerer som samlende applikationskontekst. Den holder styr på, om brugeren er autentificeret, hvilken profil og plan der er aktiv, hvilken træningssession der er i gang, og hvilke tilstande der skal være synlige for brugerfladen. Det er særligt vigtigt i Træningsmester, fordi brugerens næste relevante handling afhænger af både planlagte data og faktisk adfærd.

Dataadgang samles i repositories og services. Skærme kalder ikke databasen direkte. Laget samler læsning, skrivning og fejlhåndtering, så completion, trackerlogs, delte programmer og cyklusdata får tydelige dataflows.

Tabel~\ref{tab:tm-container-architecture} viser denne rollefordeling, og Tabel~\ref{tab:tm-sql-fields} viser centrale felter, der understøtter blandt andet idempotency, completion, aktive sessioner og cyklusprogression. Disse felter er arkitektonisk vigtige, fordi datakvalitet ikke kun handler om at have data, men om at data kan fortolkes konsistent på tværs af brugerhandlinger og systemflader \cite{mohammed2025_five_facets_data_quality}.

\subsection{Backend som sikkerhedsgrænse}

Backendens rolle er at være den autoritative sikkerheds- og datagrænse. Klienten kan give brugeren en struktureret oplevelse, men den må ikke alene afgøre, hvilke data brugeren har adgang til, eller hvilke privilegerede handlinger der er tilladte. Dette følger både mobile sikkerhedsprincipper og Supabase-modellen, hvor RLS og API-beskyttelse skal begrænse adgang efter bruger, rolle og kontekst \cite{owasp2024_masvs,supabase_rls_docs,supabase_securing_api_docs}.

Figur~\ref{fig:tm-security-boundary-summary} viser grænsen mellem mobil klient, RLS og Edge Functions. RLS beskytter brugerbundne læse- og skriveoperationer, mens Edge Functions afgrænser handlinger, der kræver server-side vurdering, integrationer eller privilegerede nøgler. Tabel~\ref{tab:tm-authorization-dcr} beskriver autorisationsreglerne som proceslogik, og Tabel~\ref{tab:tm-rls-matrix} viser, hvordan centrale datadomæner kobles til klientbeskyttelse og serverbeskyttelse.

Denne arkitektur reducerer risikoen for, at en fejl i brugerfladen giver bredere adgang end tilsigtet. Den dokumenterer dog ikke fuld sikkerhed i sig selv. RLS, helper functions, RPC og Edge Functions skal fortsat vurderes gennem test, deployment-review og konkrete adgangsscenarier.

\subsection{watchOS og Live Activity som friktionsreducerende arkitektur}

watchOS og Live Activity er arkitektoniske greb til at reducere friktion under træning. De er ikke selvstændige hovedformål, men støtteflader, der skal gøre det lettere at følge, logge eller afslutte en træning, når brugeren er midt i et træningspas. Apple beskriver Live Activities som en måde at vise løbende, tidsfølsom status, mens Watch Connectivity kan bruges til at koordinere data mellem iPhone og Apple Watch \cite{apple_activitykit_live_data_docs,apple_hig_live_activities,apple_watchconnectivity_docs}.

Figur~\ref{fig:tm-watch-sync-sequence} viser, hvordan watchOS-sporet relaterer sig til auth-sync, token-cache og RLS-baserede workoutdata. Figur~\ref{fig:tm-live-activity-idempotency-sequence} viser, hvordan Live Activity-sporet kobles til intent-write, App Group snapshot og idempotent trackerlog. Begge spor afhænger af, at appens state og backendkontrakter er fælles nok til, at flere flader kan arbejde med samme træningsforløb.

Det arkitektoniske tradeoff er, at flere flader giver bedre tilgængelighed, men også øger behovet for konsistent state, idempotente writes og tydelige sikkerhedsgrænser. Derfor er watchOS og Live Activity placeret som forlængelser af den samme træningsmodel, ikke som separate datamodeller.

\subsection{Arkitekturens afgrænsning og overgang til interaktionsdesign}

Arkitekturen forklarer, hvordan Træningsmester kan forbinde klient, backend, platformservices og sikkerhedsregler. Den forklarer derimod ikke alene, om brugeroplevelsen føles motiverende, eller om feedbacken faktisk hjælper brugeren med at fastholde træning. Det kræver en analyse af de konkrete interaktionsflows, skærme og feedbackmekanismer.

Figur~\ref{fig:tm-import-ai-edge-boundary} viser desuden, at import- og AI-sporet ligger bag backendgrænsen. Eksterne integrationer kan støtte programimport og automatisering, men de skal følge samme principper om adgangskontrol, datakvalitet og brugerbundet kontekst.

Samlet viser kapitlet, at Træningsmesters arkitektur er bygget til at holde plan, udførelse, completion, progression og sikker dataadgang sammen. Det næste spørgsmål er, hvordan denne tekniske struktur omsættes til konkrete brugerflader og handlinger. Det leder videre til interaktionsdesign.
